\documentclass[11pt,a4paper]{article}
\usepackage[T1]{fontenc}
\usepackage[utf8]{inputenc}
\usepackage{lmodern}
\usepackage[a4paper,margin=2.6cm]{geometry}
\usepackage{microtype}
\usepackage{amsmath,amssymb}
\usepackage{booktabs}
\usepackage{enumitem}
\usepackage{xcolor}
\usepackage{hyperref}
\hypersetup{colorlinks=true,linkcolor=black,citecolor=black,urlcolor=blue}
\setlist{nosep}
\setlength{\parskip}{0.6em}
\setlength{\parindent}{0pt}

\title{\textbf{Freedom of Conscience and the Right to a Correctable Process}\\\large A Candidate Symmetry for Human Rights in Interdependent Life}
\author{Albert Jan van Hoek}
\date{September 2026}

\begin{document}
\maketitle

\begin{abstract}
Modern human-rights law strongly protects the internal freedom to think, believe, and hold opinions without coercion. That protection is indispensable. Yet interdependent life creates a complementary problem: a person's private certainty can be translated into decisions that bind other people. The logical argument developed in \emph{The Room} shows why certainty alone cannot function as a truth guarantee and why a learning system loses robust correctability when potentially corrective routes are made permanently inaccessible. This paper asks how that logic might be translated into rights language. It does not propose a duty to doubt one's religion, conscience, or opinion. Instead, it develops a candidate right to procedural corrigibility: where shared decisions materially affect a person, no participant's conviction, status, or certainty should by itself create unreviewable authority, and restrictions on corrective participation should themselves remain reasoned, challengeable, and revisable. The result creates a symmetry between protection from coercion of conscience and protection from uncorrectable domination in shared processes.
\end{abstract}

\section{The existing architecture}

The Universal Declaration of Human Rights (UDHR) protects freedom of thought, conscience, religion, and opinion. Article 18 protects freedom of thought, conscience, and religion, including freedom to change religion or belief. Article 19 protects freedom of opinion and expression. Article 28 recognizes an entitlement to a social and international order in which the Declaration's rights can be realized, while Article 29 recognizes duties to the community and permits lawful limitations needed to respect the rights and freedoms of others and the general welfare of a democratic society. Article 30 prevents the Declaration from being interpreted as conferring a right to destroy the rights it protects \cite{udhr}.

The International Covenant on Civil and Political Rights (ICCPR) gives the internal dimension additional legal force. Article 18 protects freedom of thought, conscience, and religion and prohibits coercion that would impair the freedom to have or adopt a religion or belief. Manifestation of religion or belief can be limited under the conditions specified in Article 18(3) to protect public interests and the fundamental rights and freedoms of others. Article 19 protects the right to hold opinions without interference and the freedom to seek, receive, and impart information, subject to the Covenant's specified restrictions \cite{iccpr}. The Human Rights Committee has emphasized both the breadth of Article 18 and the importance of free opinion and expression for transparency, accountability, and democratic society \cite{gc22,gc34}.

European human-rights doctrine uses a similar distinction. Article 9 of the European Convention on Human Rights protects freedom of thought, conscience, and religion. The European Court of Human Rights treats the internal freedom to hold or change belief as absolute and unqualified, while manifestations may be limited under Article 9(2) because their exercise can affect other people \cite{echr9guide}.

The existing framework therefore already recognizes a basic asymmetry between the inner and outer domains:
\[
\boxed{\text{the inner conviction is strongly protected}}
\qquad\text{while}\qquad
\boxed{\text{its outward effects can be legally constrained}.}
\]

This paper does not challenge that architecture. It asks whether an additional general procedural principle can make its relational logic more explicit.

\section{The gap exposed by The Room}

Human-rights instruments contain many procedural protections in particular contexts: equality before courts, fair hearings, effective remedies, review mechanisms, political participation, and protections for receiving and imparting information. The ICCPR, for example, requires accessible and effective remedies for Covenant violations and protects fair and public hearings in appropriate legal proceedings \cite{iccpr,gc31}.

The instruments examined here, however, do not express as a single general right the following learning-system principle:
\begin{quote}
\textbf{When a shared decision materially affects a person, the decision process should remain capable of receiving relevant correction, including correction directed at its own restrictions on participation.}
\end{quote}

That proposed principle does not arise because every person is equally likely to be right. It arises because the system generally lacks an infallible truth oracle. If several sincere agents can hold incompatible beliefs, certainty alone cannot identify the true belief. If a participant is affected by a decision, and the current system can make that participant's corrective input categorically inadmissible, then the system can enter a state in which it is wrong and has disabled one route by which the error could be exposed.

The legal question is therefore not whether the state may force citizens to become psychologically uncertain. It should not. The question is whether private certainty may become \emph{unreviewable shared authority}.

\section{A proposed symmetry}

The symmetry can be expressed in two directions:
\[
\boxed{\text{No one may force my conviction.}}
\]
\[
\boxed{\text{My conviction alone may not make our shared process uncorrectable.}}
\]

Or, in institutional language:
\begin{center}
\begin{tabular}{p{0.43\textwidth}p{0.43\textwidth}}
\toprule
\textbf{Freedom of conscience} & \textbf{Procedural corrigibility}\\
\midrule
Protects the internal model & Protects the corrective relationship\\
No forced belief & No unreviewable shared authority\\
Freedom to remain convinced & Route to challenge consequential decisions\\
Protection from epistemic coercion & Protection from epistemic domination\\
Plurality of models & Correctable interaction among models\\
\bottomrule
\end{tabular}
\end{center}

This is a structural symmetry, not a claim that the two protections would have identical doctrinal status, limitation tests, or legal weight. Existing law deliberately gives particularly strong protection to the internal forum, while procedural rights often depend on context. The point is different: a rights architecture for interdependent fallible agents may need protection in both directions.

\section{Candidate right to a correctable process}

The following is draft language for discussion, not proposed treaty text ready for adoption.

\begin{quote}
\textbf{Article X -- Right to a Correctable Process}

\textbf{1.} Everyone whose rights, freedoms, or legitimate interests are materially affected by a shared public or institutional decision is entitled, on equal terms, to a process capable of receiving relevant reasons, evidence, objections, and corrections.

\textbf{2.} No person's belief, opinion, conviction, status, or certainty shall by itself confer unreviewable authority to determine a shared outcome affecting others.

\textbf{3.} Decisions may be settled and acted upon. Where material grounds arise to question their factual basis, reasoning, fairness, or continued applicability, an effective route to challenge, review, or revision shall remain available under rules appropriate to the context.

\textbf{4.} A restriction on a person's access to such a process may be imposed where lawful and necessary, but the restriction shall be reasoned, proportionate to its purpose, and itself subject to an appropriate route of challenge or review. No restriction shall become immune from correction merely by virtue of the authority that imposed it.

\textbf{5.} Nothing in this Article requires any person to abandon, modify, disclose, or doubt a privately held thought, conscience, religion, belief, or opinion. Nor does it require equal weight to be given to claims, evidence, expertise, or conduct that are not equal in relevance or quality.
\end{quote}

The fifth paragraph is essential. The proposal does not turn freedom of conscience into a legal obligation to say ``I may be wrong.'' It protects the person's internal freedom while constraining the institutional inference from certainty to unilateral authority.

\section{Why this is not a right to endless debate}

A correctable process must still be capable of deciding. Human systems have finite time and attention. Courts need finality. Public agencies need to act. Scientific institutions reject weak claims. Employers and schools impose boundaries. Dangerous or abusive behavior may justify restrictions.

Procedural corrigibility therefore cannot mean permanent openness. It is better captured by the principle
\[
\boxed{\text{no uncorrectable closure}.}
\]

A final decision can remain final under ordinary circumstances while still admitting exceptional reopening for genuinely new evidence, material error, procedural unfairness, changed circumstances, or another recognized ground. An excluded participant need not retain unrestricted access, but the exclusion can be made conditional, reviewable, time-limited, independently appealable, or capable of reversal when the reasons for it no longer hold.

Law already uses many such devices. Appeals, judicial review, review of administrative decisions, reopening on new evidence, requirements to give reasons, proportionality review, and effective remedies are all mechanisms by which legal systems prevent some decisions from becoming completely self-sealing. The proposed right does not replace these doctrines. It offers a general principle that can explain why they matter as a family: they preserve the capacity of the institutional system to correct consequential error.

\section{A right over process rather than conclusion}

The proposal is intentionally agnostic about who is correct. It protects neither Alice's conclusion nor Bob's conclusion. It protects the architecture through which either conclusion can be challenged when it governs shared life.

This is especially important where the decision-maker believes the excluded person to be unreliable. A rule such as ``nothing from B may ever matter'' evaluates future evidence solely by source identity. It does not merely assign lower weight; it prevents the system from discovering cases in which its own source judgement was mistaken.

A legal right built around correctability would therefore protect \emph{standing and routes}, not epistemic equality. It would be compatible with:
\begin{itemize}
  \item differential expertise;
  \item evidentiary thresholds;
  \item sanctions for abuse;
  \item deadlines and finality rules;
  \item filtering of irrelevant or repetitive material;
  \item emergency decisions followed by later review.
\end{itemize}

Its core concern would be whether the architecture can turn a current judgement into an institutionally unrevisable judgement about who may ever count as a corrective source.

\section{From human rights to a learning constitution}

The candidate Article X can also be understood as a constitutional property of a learning system. A society cannot guarantee that its decisions will be true. It can, however, choose rules that reduce the number of states in which a current error is protected from every admissible correction.

That reframes several familiar rights. Freedom of expression is not only valuable because speaking is intrinsically important; it can preserve information pathways. Fair-hearing rights are not only dignitary protections; they can expose errors before coercive decisions become final. Effective remedies do not only compensate for wrongs; they keep the system capable of revising itself. Equal procedural standing prevents current authority from converting itself into a permanent truth privilege.

The point is not to reduce existing rights to an instrumental theory of learning. Human dignity is not merely an optimization objective. Rather, the learning-system perspective reveals an additional structural reason why some rights fit together: fallible agents need protected pathways through which consequential shared models can still be challenged.

\section{What the formal logic can support}

The accompanying Lean model supports conditional statements such as:
\begin{itemize}
  \item pairwise incompatible claims cannot all be true;
  \item incompatible certainties cannot both be truth-guaranteed;
  \item incompatible exclusive requirements cannot share one mutually acceptable outcome;
  \item permanent exclusion of a potentially corrective source defeats a defined form of source-robust correctability;
  \item a learning constitution that requires restrictions to remain challengeable and revisable excludes self-sealing restrictions;
  \item if that constitution is preserved by every allowed transition, the property remains true in every reachable state.
\end{itemize}

Those are deductive results under explicit definitions and premises. They do not establish the normative statement that international law ought to adopt Article X. The normative bridge remains a human choice: whether societies wish to protect equal standing, non-domination, and continuing institutional correctability strongly enough to constitutionalize them.

This separation is a strength. The legal proposal can be debated without pretending that Lean proves morality, and the formal argument can be inspected without requiring agreement with the legal proposal.

\section{The resulting perspective}

The combined architecture can be stated very simply:
\begin{quote}
\textbf{Nobody may own your mind. Nobody's mind should become uncorrectable authority over you.}
\end{quote}

The first sentence captures the protective insight already central to freedom of thought, conscience, religion, and opinion. The second is the proposed complement for interdependent life.

The symmetry does not require citizens to converge on one worldview. On the contrary, it permits deep and persistent disagreement. What it protects is the possibility that disagreement can still carry corrective information into shared institutions.

That may be the more realistic foundation for peaceful pluralism among fallible agents: not agreement, not compulsory doubt, but a constitution that prevents present certainty from permanently abolishing future correction.

\section*{Status and next steps}

This is a conceptual and legal-design proposal, not a statement of existing law and not legal advice. The instruments cited below were used to identify the existing protections against coercion of belief, the qualified nature of manifestations and expression, and existing procedural/remedial protections. A serious legal-development programme would require comparative constitutional analysis, jurisprudential review, institutional design, proportionality safeguards, disability and child-rights analysis, democratic-process analysis, and scrutiny of interactions with due process, finality, privacy, safety, equality, and freedom of association.

The immediate contribution is narrower: the logic identifies a problem, and the draft Article X provides one possible legal language for addressing it without weakening freedom of conscience.

\begin{thebibliography}{9}

\bibitem{udhr}
United Nations General Assembly.
\emph{Universal Declaration of Human Rights}, Resolution 217 A (III), 10 December 1948, especially Articles 18, 19, 28--30.
\url{https://www.ohchr.org/en/human-rights/universal-declaration/translations/english}

\bibitem{iccpr}
United Nations General Assembly.
\emph{International Covenant on Civil and Political Rights}, Resolution 2200A (XXI), 16 December 1966, especially Articles 2(3), 14, 18 and 19.
\url{https://www.ohchr.org/sites/default/files/ccpr.pdf}

\bibitem{gc22}
UN Human Rights Committee.
\emph{General Comment No. 22: Article 18 (Freedom of Thought, Conscience or Religion)}, CCPR/C/21/Rev.1/Add.4, 1993.
\url{https://docstore.ohchr.org/}

\bibitem{gc34}
UN Human Rights Committee.
\emph{General Comment No. 34: Article 19: Freedoms of Opinion and Expression}, CCPR/C/GC/34, 12 September 2011.
\url{https://docstore.ohchr.org/}

\bibitem{gc31}
UN Human Rights Committee.
\emph{General Comment No. 31: The Nature of the General Legal Obligation Imposed on States Parties to the Covenant}, CCPR/C/21/Rev.1/Add.13, 2004.
\url{https://docstore.ohchr.org/}

\bibitem{echr9guide}
European Court of Human Rights.
\emph{Guide on Article 9 of the European Convention on Human Rights: Freedom of Thought, Conscience and Religion}, updated 5 March 2026.
\url{https://ks.echr.coe.int/documents/d/echr-ks/guide_art_9_eng-pdf}

\end{thebibliography}

\end{document}
